\documentclass{article}

\usepackage[utf8]{inputenc}
\usepackage[T1]{fontenc}
\usepackage{amsmath}
\usepackage{amssymb}
\usepackage{booktabs}
\usepackage{graphicx}
\usepackage{multirow}
\usepackage{xcolor}
\usepackage{url}
\usepackage{hyperref}
\usepackage[capitalize]{cleveref}
\usepackage{microtype}
\usepackage{geometry}
\geometry{margin=1in}

\title{Latent Bridge Matching for Fast Multi-Style Artistic Transfer}
\author{Anonymous Authors}
\date{}

\begin{document}
\maketitle

\begin{abstract}
TODO: Abstract.
\end{abstract}

\section{Introduction}
Artistic style transfer has long served as a testbed for separating image content from visual appearance. Classical neural style transfer showed that convolutional features and Gram statistics can synthesize compelling artistic textures \cite{gatys2016image}, while feed-forward arbitrary style transfer made real-time deployment practical through normalization-based style modulation \cite{huang2017adain}. More recent attention, transformer, training-free, and diffusion-based systems have improved flexibility and style fidelity \cite{park2019sanet,deng2022stytr2,hong2023aespa,huang2024aesfa,gim2024cast,chung2024styleid}. However, practical multi-style artistic transfer still exposes a persistent three-way tension: strong target-style expression, faithful content preservation, and low computational cost.

This tension is especially visible in domain-level multi-style transfer. Efficient multi-style systems can support many styles and fast inference, but the output may become washed, over-smoothed, or contaminated by structured grain artifacts. Diffusion-based or training-free methods can produce strong style semantics, but may be expensive or suffer semantic drift under content-preserving evaluation. In our local strict-750 protocol, for example, StyleID obtains the highest raw CLIP-style score but exhibits severe content collapse, while SaMST obtains strong raw style and content scores but produces visible muddy and grain-like texture artifacts. These observations suggest that standard style/content metrics alone are insufficient: a competitive system must also be evaluated for perceptual quality and artifact-sensitive texture realism.

We propose a fast latent-space multi-style artistic transfer framework that models stylization as a style-conditioned latent probability flow. Instead of directly predicting RGB images, the method transports a compact VAE latent through a neural velocity field. A terminal sliced-Wasserstein distance (SWD) term aligns generated endpoints with the target artistic domain, while kinetic regularization penalizes excessive latent displacement and preserves content. A learnable style spatial prior and semantic cross-attention provide a lightweight approximation to local style transport without per-image optimization or diffusion sampling.

Our contributions are:
\begin{itemize}
    \item We formulate domain-level multi-style artistic transfer as a bridge-inspired latent transport problem with separate terminal style matching and path-energy content preservation.
    \item We introduce an efficient style spatial prior and semantic cross-attention architecture that supports fast style-conditioned latent velocity prediction.
    \item We provide a strict 750-output evaluation protocol with complete local comparisons to AdaIN, S2WAT, SaMST, StyleID, and our method, together with artifact-sensitive diagnostics for grain and perceptual quality.
    \item We conduct destructive ablations and 40-run category-weight sweeps showing that terminal SWD drives style, kinetic regularization preserves content, and naive color matching harms the style-content trade-off.
\end{itemize}

\section{Related Work}
\paragraph{Classical and arbitrary neural style transfer.}
Neural style transfer began with optimization-based feature reconstruction and Gram-matrix style losses \cite{gatys2016image}. AdaIN later enabled real-time arbitrary transfer by aligning channel-wise feature statistics \cite{huang2017adain}. Attention-based methods such as SANet improved local style matching through style-attentional feature aggregation \cite{park2019sanet}. Recent aesthetic and pattern-aware methods, including AesPA-Net and AesFA, further emphasize pattern repeatability, aesthetic features, and high-quality arbitrary stylization \cite{hong2023aespa,huang2024aesfa}. These methods demonstrate the value of feature statistics and local matching, but they generally do not explicitly model a style-conditioned transport path with a terminal distribution constraint.

\paragraph{Transformer and multi-style transfer.}
Transformer-based style transfer methods such as StyTr2 and S2WAT use attention mechanisms to improve content-style interaction and long-range feature modeling \cite{deng2022stytr2,xia2024s2wat}. Multi-style approaches aim to amortize many styles in a single compact system; SaMST is a recent representative emphasizing pluggable style representations, low parameter count, and fast inference \cite{liu2024samst}. These systems are directly related to our goal of efficient multi-style transfer. In contrast to representation-only style injection, our method embeds semantic cross-attention inside a latent bridge objective, allowing terminal distribution matching and kinetic regularization to explicitly shape the style-content trade-off.

\paragraph{Training-free and diffusion-based style transfer.}
Training-free and diffusion-based approaches provide a different route to flexible style transfer. CAST uses VQ autoencoders for content-adaptive training-free style transfer \cite{gim2024cast}, while StyleID injects style into large-scale diffusion models without training \cite{chung2024styleid}. These methods can be visually strong, but diffusion sampling and semantic drift can limit their suitability for fast domain-level transfer. Our work targets a complementary efficiency regime: compact latent integration with style-id conditioning rather than iterative diffusion sampling or per-style optimization.

\paragraph{Distribution matching and perceptual evaluation.}
Our objective is motivated by optimal transport and distribution matching. Sinkhorn distances and computational optimal transport provide efficient tools for approximate transport problems \cite{cuturi2013sinkhorn,peyre2019computational}, while the Schrodinger problem connects entropic transport to stochastic path measures \cite{leonard2014survey}. We use SWD as a lightweight terminal distribution matching objective. For evaluation, we combine perceptual and distributional measures including LPIPS \cite{zhang2018lpips}, FID \cite{heusel2017fid}, KID \cite{binkowski2018kid}, DISTS \cite{ding2020dists}, MUSIQ \cite{ke2021musiq}, MANIQA \cite{yang2022maniqa}, and style-transfer-specific ArtFID \cite{wright2022artfid}. This broad metric set is important because raw CLIP-style and structural scores may fail to penalize grain-like artifacts.

\section{Method}
TODO: Method. To be filled by the section-writing agent.

\section{Experiments}
TODO: Experiments. To be filled by the section-writing agent.

\section{Discussion and Limitations}
TODO: Discussion and limitations. To be filled by the section-writing agent.

\section{Conclusion}
TODO: Conclusion. To be filled by the section-writing agent.

\bibliographystyle{plain}
\bibliography{refs}

\end{document}

